\section{A constructive interface for height descent and finite generation}
\label{sec:generation}

An important role of height in Diophantine geometry is to reduce all points
to finitely many small ones. The following theorem isolates a version of
that argument in which the required witnesses are supplied by actual
primitive recursive functions. It is a theorem about explicitly encoded
groups, not an assertion that every elliptic curve already comes equipped
with these data.

\begin{theorem}[Finite generation from a computable reducing step]
\label{thm:generation}
Fix an abelian group represented by natural-number codes, with primitive
recursive operations, equality and a validity predicate. Suppose $\iPRA$
proves its group identities on valid codes. Let $h$ be a primitive recursive
natural-valued height. Fix $m\ge2$, a threshold $T$, and finite lists of valid group elements
\[
      R=(r_0,\ldots,r_{s-1}),\qquad S=(s_0,\ldots,s_{t-1}).
\]
Suppose the following data and their correctness proofs are given:
\begin{enumerate}
\item $S$ lists all valid group elements of height at most $T$;
\item primitive recursive functions $D(P)$ and $i(P)$ such that, for every
valid $P$ with $h(P)>T$, the point $D(P)$ is valid, $i(P)<s$, and
\[
               P=mD(P)+r_{i(P)},\qquad h(D(P))<h(P).
\]
\end{enumerate}
Then $R\cup S$ generates the group. More precisely, a primitive recursive
algorithm produces, for each valid $P$, explicit integer coefficients
expressing $P$ in these generators. Its correctness is provable in $\iPRA$.
\end{theorem}
\begin{proof}
Keep a current point, an active-step count $k$, and a coefficient vector
$(a_0,\ldots,a_{s-1})$. Start at $P$, $k=0$, and all coefficients zero.
When the current point $Q$ has height at most $T$, keep the whole state
unchanged. Otherwise first save $i=i(Q)$ and $D(Q)$ from the old point,
then replace $Q$ by $D(Q)$, add $m^k$ to $a_i$, and increase $k$ by one.
Run this stopped iteration for exactly
$h(P)\mathbin{\dotminus}T$ stages, where $\dotminus$ is truncated
subtraction.

The rank $h(Q)\mathbin{\dotminus}T$ is zero precisely in the stopped
case. In an active step it decreases by at least one; in a stopped step
it remains zero. Quantifier-free induction therefore proves that after
$j$ stages it is at most
$(h(P)\mathbin{\dotminus}T)\mathbin{\dotminus}j$. The final point has
height at most $T$ and hence is some $s_\ell$ in the supplied list. Its
index is found by bounded equality search.

Throughout the iteration the following group identity is maintained:
\begin{equation}\label{eq:descent-expansion}
                   P=m^k Q+\sum_{i<s}a_i r_i.
\end{equation}
It holds initially. An active step substitutes
$Q=mD(Q)+r_{i(Q)}$ into the right side and uses distributivity of
integer multiplication in an abelian group. The result is exactly the
coefficient update. A stopped step changes nothing. At the final state,
replace $Q$ by $s_\ell$ in \eqref{eq:descent-expansion}; this is the
promised expression. All coefficients can even be chosen nonnegative
with this particular supplied reduction.

The loops have the given numerical fuel. Scalar multiplication is repeated
group addition, so its definition is primitive recursive. Every exponent
is at most $h(P)\dotminus T$, and every coefficient is bounded by
$\sum_{j\le h(P)}m^j$. Equation \eqref{eq:descent-expansion} and the rank
inequality are quantifier-free statements about the coded operations, so
their inductions are instances of the permitted rule. No choices from an
unspecified finite quotient are made.
\end{proof}

A commonly useful sufficient numerical condition for strict decrease is
\[
              2h(D(P))\le h(P)+T\qquad(h(P)>T).
\]
Indeed $h(P)+T<2h(P)$ implies $h(D(P))<h(P)$ by integer arithmetic.
This avoids real logarithms and limit-defined canonical heights in the
interface. Constructing such a height and a reducing selector for a
particular group is an additional mathematical obligation.

\subsection{What this does and does not establish about larger theorems}

Theorem~\ref{thm:generation} captures the final reduction step of a
Mordell--Weil-style proof. A statement that a quotient group is finite is
not, on its own, the supplied primitive recursive selector required here.
One must provide representatives, produce the corresponding preimages,
and prove the height decrease in the chosen arithmetic. Likewise a proof
of finiteness of rational points is not automatically an explicit global
height bound to which Corollary~\ref{cor:heightbound} applies.

The results actually completed in this paper need none of those additional
inputs: pointed conics come with explicit formulas, rational preperiodic
points have the resultant-derived height bound, and the curve $E_1$ is
settled by its concrete decreasing transformation. The broad classical
theorems are valuable targets for further work, but they are not hidden
premises in these arguments. No inference here rests on an assertion that
an unexamined classical proof must formalize in weak arithmetic.

\subsection{Logical inventory of the completed proofs}

The paper uses first-order numerical quantification to state its theorems,
including unbounded universal claims about all rational points. Their
proofs use the following finite mechanisms:
\begin{center}
\begin{tabular}{p{0.31\textwidth}p{0.60\textwidth}}
\toprule
Mathematical operation & Arithmetic implementation \\
\midrule
Projective equivalence & Gcd normalization and Bezout uniqueness \\
Finite point loci & Explicit coordinate boxes and polynomial tests \\
Coordinate morphisms & Integral homogeneous formulas and nonvanishing guards \\
Height lower bounds & Polynomial certificates and bounds on their common divisors \\
Infinite families & An injective formula evaluated on a supplied finite initial segment \\
Preperiodicity & A proved height cutoff and a finite repeated-state search \\
Nonexistence by descent & Iteration for the initial integer rank \\
Local impossibility & Exhaustive residue tests for a supplied modulus \\
Finite generation & A supplied reducing selector and a telescoping group identity \\
\bottomrule
\end{tabular}
\end{center}
Each induction has an equation, a Boolean invariant, or a bounded
arithmetical invariant. There is no quantification over arbitrary sets of
natural numbers, no completeness principle for $\R$, and no appeal to an
unspecified unbounded search. This is the substantive reason the development
fits the chosen constructive arithmetic.
